% --- Archon physics formalization source begin ---
% archon:physics
% archon:covers PhyXMiniProblems/problem_phyx_mini_0604.lean
% archon:source-report reports/phyx_mini/problem_phyx_mini_0604.source.json

\chapter{Physics problem phyx\_mini\_0604}
\label{ch:PhyXMiniProblems_problem_phyx_mini_0604}

\paragraph{Problem source.}
\#\# Physical scenario

The Boltzmann equation \$\$P(n) = \textbackslash{}frac\{1\}\{Z\} e\textasciicircum{}\{-eta E\_n\}, \textbackslash{}quad Z = \textbackslash{}sum\_n e\textasciicircum{}\{-eta E\_n\}, \textbackslash{}quad eta = \textbackslash{}frac\{1\}\{k\_B T\},\$\$ gives the probability of finding a system in the state \$n\$ (with energy \$E\_n\$), at temperature \$T\$ (\$k\_B\$ is Boltzmann's constant). Note: The probability here refers to the random thermal distribution, and has nothing to do with quantum indeterminacy. Quantum mechanics will only enter this problem through quantization of the energies \$E\_n\$.

\#\# Question

A crystal consisting of \$N\$ atoms can be thought of as a collection of \$3N\$ oscillators (each atom is attached by springs to its 6 nearest neighbors, along the \$x, y\$, and \$z\$ directions, but those springs are shared by the atoms at the two ends). The heat capacity of the crystal (per atom) will therefore be \$\$C = 3 \textbackslash{}frac\{dar\{E\}\}\{dT\}.\$\$ Show that (in this model) \$\$C = 3k\_B \textbackslash{}left( \textbackslash{}frac\{\textbackslash{}Theta\_E\}\{T\} 
ight)\textasciicircum{}2 \textbackslash{}frac\{e\textasciicircum{}\{\textbackslash{}Theta\_E/T\}\}\{(e\textasciicircum{}\{\textbackslash{}Theta\_E/T\} - 1)\textasciicircum{}2\},\$\$ where \$\textbackslash{}Theta\_E \textbackslash{}equiv \textbackslash{}hbar \textbackslash{}omega / k\_B\$ is the so-called Einstein temperature. The same reasoning using the classical expression for \$ar\{E\}\$ yields \$C\_\{\textbackslash{}text\{classical\}\} = 3k\_B\$, independent of temperature.

\#\# Answer choices

- A: A: 3k\_B \textbackslash{}cdot \textbackslash{}frac\{e\textasciicircum{}\{\textbackslash{}Theta\_E/T\}\}\{(e\textasciicircum{}\{\textbackslash{}Theta\_E/T\} - 1)\}

- B: B: 3k\_B \textbackslash{}cdot \textbackslash{}left( \textbackslash{}frac\{\textbackslash{}Theta\_E\}\{T\} 
ight)\textasciicircum{}2 \textbackslash{}cdot \textbackslash{}frac\{1\}\{e\textasciicircum{}\{\textbackslash{}Theta\_E/T\} - 1\}

- C: C: 3k\_B \textbackslash{}left( \textbackslash{}frac\{\textbackslash{}Theta\_E\}\{T\} 
ight)\textasciicircum{}2 \textbackslash{}frac\{e\textasciicircum{}\{\textbackslash{}Theta\_E/T\}\}\{(e\textasciicircum{}\{\textbackslash{}Theta\_E/T\} - 1)\textasciicircum{}2\}

- D: D: 3k\_B

\#\# Figure caption (auxiliary; use the image as primary evidence)

The image is a graph showing the specific heat (C) as a function of temperature (T) in Kelvin (K). 

- **Axes:** 
  - The x-axis represents temperature, ranging from 0 K to 1000 K.
  - The y-axis represents specific heat, in units of J/g·K, ranging from 0 to 2.0.

- **Data Points:** 
  - There are two sets of data points represented as curves:
    - \textbackslash{}( C\_v \textbackslash{}) is denoted by solid black circles.
    - \textbackslash{}( C\_p \textbackslash{}) is denoted by open black circles.

- **Trends:**
  - Both curves show an increase in specific heat with increasing temperature.
  - \textbackslash{}( C\_p \textbackslash{}) is slightly higher than \textbackslash{}( C\_v \textbackslash{}) as temperature increases.

- **Labels:**
  - The curves are labeled \textbackslash{}( C\_v \textbackslash{}) and \textbackslash{}( C\_p \textbackslash{}) near the top right of the graph.

The graph plots specific heat against temperature, showing how specific heat values increase with temperature for the two conditions.

\#\# Dataset metadata

Quantum Phenomena; Numerical Reasoning, Physical Model Grounding Reasoning

\paragraph{Recorded answer/context.}
C: C: 3k\_B \textbackslash{}left( \textbackslash{}frac\{\textbackslash{}Theta\_E\}\{T\} 
ight)\textasciicircum{}2 \textbackslash{}frac\{e\textasciicircum{}\{\textbackslash{}Theta\_E/T\}\}\{(e\textasciicircum{}\{\textbackslash{}Theta\_E/T\} - 1)\textasciicircum{}2\}

\paragraph{Figure/image path.}
/root/proposal\_for\_physic/science-mango/phyx\_mini\_run/phyx\_data/test\_image/604.png

\paragraph{Formalization target.}
create a compiling Lean file with sorry bodies at `PhyXMiniProblems/problem\_phyx\_mini\_0604.lean`. The Lean declarations must preserve the physical quantities, dimensions or dimensional roles, figure labels, governing-law hypotheses, and final relation expressed by this problem.
Use Mathlib/Physlib names found through LeanExplore where available. If a physics API is missing, introduce faithful local abstractions rather than scalar placeholder aliases.

\begin{theorem}[Physics formalization target]
\label{thm:physics:phyx_mini_0604:target}
\lean{PhyXMini0604.classical_heat_capacity}
\uses{def:physics:phyx-mini-0604:phyxmini0604-classicalheatcapacityperatom}
The assigned autoformalize agent should translate this physics problem into Lean declarations in the covered file, with theorem and lemma proof bodies written as `by sorry`.
\end{theorem}
\begin{proof}
This is an autoformalization task, not a proof task. Produce statements that can later be proved by the physics prover without weakening the physical model.
\end{proof}

% --- Archon declaration topology begin ---
\section{Declaration topology}
\begin{definition}[Cartesian Axis]
  \label{def:physics:phyx-mini-0604:phyxmini0604-cartesianaxis}
  \lean{PhyXMini0604.CartesianAxis}
  The three Cartesian directions explicitly named in the crystal model
\end{definition}
\begin{proof}
  This is the declared model definition of Cartesian Axis.
\end{proof}
\begin{definition}[Axis Orientation]
  \label{def:physics:phyx-mini-0604:phyxmini0604-axisorientation}
  \lean{PhyXMini0604.AxisOrientation}
  The two nearest-neighbor orientations along any Cartesian axis
\end{definition}
\begin{proof}
  This is the declared model definition of Axis Orientation.
\end{proof}
\begin{definition}[Nearest Neighbor Direction]
  \label{def:physics:phyx-mini-0604:phyxmini0604-nearestneighbordirection}
  \lean{PhyXMini0604.NearestNeighborDirection}
  \uses{def:physics:phyx-mini-0604:phyxmini0604-cartesianaxis, def:physics:phyx-mini-0604:phyxmini0604-axisorientation}
  A nearest-neighbor direction is one of ±x, ±y, or ±z
\end{definition}
\begin{proof}
  This is the declared model definition of Nearest Neighbor Direction.
\end{proof}
\begin{theorem}[nearest Neighbor Direction card]
  \label{thm:physics:phyx-mini-0604:phyxmini0604-nearestneighbordirection-card}
  \lean{PhyXMini0604.nearestNeighborDirection_card}
  \uses{def:physics:phyx-mini-0604:phyxmini0604-nearestneighbordirection}
  There are six nearest-neighbor directions in the stated crystal geometry
\end{theorem}
\begin{proof}
  The conclusion follows from the governing relations and figure readouts named in the statement of nearest Neighbor Direction card.
\end{proof}
\begin{definition}[Crystal Geometry]
  \label{def:physics:phyx-mini-0604:phyxmini0604-crystalgeometry}
  \lean{PhyXMini0604.CrystalGeometry}
  Aggregate geometric data for the spring model of the crystal. The equality oscillatorCount = 3 * atomCount records the sharing of each of the six nearest-neighbor springs by the two atoms at its endpoints
\end{definition}
\begin{proof}
  This is the declared model definition of Crystal Geometry.
\end{proof}
\begin{definition}[Einstein Crystal Model]
  \label{def:physics:phyx-mini-0604:phyxmini0604-einsteincrystalmodel}
  \lean{PhyXMini0604.EinsteinCrystalModel}
  \uses{def:physics:phyx-mini-0604:phyxmini0604-crystalgeometry}
  The Einstein model consists of identical one-dimensional quantum oscillators and their discrete canonical ensemble. The counting-measure and zero-degree-of-freedom fields express that ℕ labels discrete quantum energy levels rather than a classical phase space. The current heat-capacity conclusion is deliberately not a field of this model
\end{definition}
\begin{proof}
  This is the declared model definition of Einstein Crystal Model.
\end{proof}
\begin{definition}[einstein Temperature]
  \label{def:physics:phyx-mini-0604:phyxmini0604-einsteintemperature}
  \lean{PhyXMini0604.einsteinTemperature}
  \uses{def:physics:phyx-mini-0604:phyxmini0604-einsteincrystalmodel}
  The Einstein temperature Θ\_E = ℏ ω / k\_B, as a temperature-valued scalar readout in the units used by Constants.kB
\end{definition}
\begin{proof}
  This is the declared model definition of einstein Temperature.
\end{proof}
\begin{definition}[heat Capacity Per Atom]
  \label{def:physics:phyx-mini-0604:phyxmini0604-heatcapacityperatom}
  \lean{PhyXMini0604.heatCapacityPerAtom}
  \uses{def:physics:phyx-mini-0604:phyxmini0604-einsteincrystalmodel}
  The heat capacity per atom in the Einstein model. Physlib's CanonicalEnsemble.heatCapacity is d Ē / dT for one oscillator, and the crystal has three such modes per atom
\end{definition}
\begin{proof}
  This is the declared model definition of heat Capacity Per Atom.
\end{proof}
\begin{theorem}[partition Function eq boltzmann sum]
  \label{thm:physics:phyx-mini-0604:phyxmini0604-partitionfunction-eq-boltzmann-sum}
  \lean{PhyXMini0604.partitionFunction_eq_boltzmann_sum}
  \uses{def:physics:phyx-mini-0604:phyxmini0604-einsteincrystalmodel}
  The canonical partition function is the sum of Boltzmann weights over the discrete oscillator levels
\end{theorem}
\begin{proof}
  The conclusion follows from the governing relations and figure readouts named in the statement of partition Function eq boltzmann sum.
\end{proof}
\begin{theorem}[probability eq boltzmann weight]
  \label{thm:physics:phyx-mini-0604:phyxmini0604-probability-eq-boltzmann-weight}
  \lean{PhyXMini0604.probability_eq_boltzmann_weight}
  \uses{def:physics:phyx-mini-0604:phyxmini0604-einsteincrystalmodel}
  The thermal probability of level n obeys the Boltzmann law. This is a classical random thermal probability; quantum mechanics enters only through the quantized energy map in EinsteinCrystalModel.energy\_quantized
\end{theorem}
\begin{proof}
  The conclusion follows from the governing relations and figure readouts named in the statement of probability eq boltzmann weight.
\end{proof}
\begin{theorem}[oscillator mean Energy eq]
  \label{thm:physics:phyx-mini-0604:phyxmini0604-oscillator-meanenergy-eq}
  \lean{PhyXMini0604.oscillator_meanEnergy_eq}
  \uses{def:physics:phyx-mini-0604:phyxmini0604-einsteincrystalmodel, def:physics:phyx-mini-0604:phyxmini0604-einsteintemperature}
  The mean energy of one quantized oscillator, including its temperature- independent zero-point energy. This is the statistical-mechanics input whose temperature derivative gives the heat capacity
\end{theorem}
\begin{proof}
  The conclusion follows from the governing relations and figure readouts named in the statement of oscillator mean Energy eq.
\end{proof}
\begin{theorem}[einstein crystal heat capacity]
  \label{thm:physics:phyx-mini-0604:phyxmini0604-einstein-crystal-heat-capacity}
  \lean{PhyXMini0604.einstein_crystal_heat_capacity}
  \uses{def:physics:phyx-mini-0604:phyxmini0604-einsteincrystalmodel, def:physics:phyx-mini-0604:phyxmini0604-einsteintemperature, def:physics:phyx-mini-0604:phyxmini0604-heatcapacityperatom}
  **** For positive absolute temperature, the heat capacity per atom of the Einstein crystal has the stated closed form
\end{theorem}
\begin{proof}
  The conclusion follows from the governing relations and figure readouts named in the statement of einstein crystal heat capacity.
\end{proof}
\begin{definition}[classical Oscillator Mean Energy]
  \label{def:physics:phyx-mini-0604:phyxmini0604-classicaloscillatormeanenergy}
  \lean{PhyXMini0604.classicalOscillatorMeanEnergy}
  The classical equipartition expression Ē = k\_B T for one oscillator, viewed as a real-valued function so that its temperature derivative is defined
\end{definition}
\begin{proof}
  This is the declared model definition of classical Oscillator Mean Energy.
\end{proof}
\begin{definition}[classical Heat Capacity Per Atom]
  \label{def:physics:phyx-mini-0604:phyxmini0604-classicalheatcapacityperatom}
  \lean{PhyXMini0604.classicalHeatCapacityPerAtom}
  \uses{def:physics:phyx-mini-0604:phyxmini0604-classicaloscillatormeanenergy}
  Classical heat capacity per atom, again with three oscillator modes per atom and with the derivative restricted to positive temperatures
\end{definition}
\begin{proof}
  This is the declared model definition of classical Heat Capacity Per Atom.
\end{proof}
\begin{definition}[Specific Heat Curve]
  \label{def:physics:phyx-mini-0604:phyxmini0604-specificheatcurve}
  \lean{PhyXMini0604.SpecificHeatCurve}
  The two curve labels visible in the auxiliary specific-heat figure
\end{definition}
\begin{proof}
  This is the declared model definition of Specific Heat Curve.
\end{proof}
\begin{definition}[Figure Marker]
  \label{def:physics:phyx-mini-0604:phyxmini0604-figuremarker}
  \lean{PhyXMini0604.FigureMarker}
  Marker styles used to distinguish the two plotted curves
\end{definition}
\begin{proof}
  This is the declared model definition of Figure Marker.
\end{proof}
\begin{definition}[Specific Heat Figure Readout]
  \label{def:physics:phyx-mini-0604:phyxmini0604-specificheatfigurereadout}
  \lean{PhyXMini0604.SpecificHeatFigureReadout}
  \uses{def:physics:phyx-mini-0604:phyxmini0604-specificheatcurve, def:physics:phyx-mini-0604:phyxmini0604-figuremarker}
  Direct readout of the auxiliary figure's axes and legend. Axis values are numerical readouts: kelvin horizontally and J g⁻¹ K⁻¹ vertically. They are metadata and are not hypotheses of the Einstein formula
\end{definition}
\begin{proof}
  This is the declared model definition of Specific Heat Figure Readout.
\end{proof}
\begin{definition}[figure604 Readout]
  \label{def:physics:phyx-mini-0604:phyxmini0604-figure604readout}
  \lean{PhyXMini0604.figure604Readout}
  \uses{def:physics:phyx-mini-0604:phyxmini0604-specificheatfigurereadout}
  Readout of figure 604. The image extends to approximately 1100 K, even though its last numbered horizontal tick is 1000 K
\end{definition}
\begin{proof}
  This is the declared model definition of figure604 Readout.
\end{proof}
% --- Archon declaration topology end ---
% --- Archon physics formalization source end ---
